\documentclass[12pt, a4paper]{article}

% ── Packages ──────────────────────────────────────────────────────────────────
\usepackage[utf8]{inputenc}
\usepackage[T1]{fontenc}
\usepackage[]{babel}
\usepackage[top=2.2cm, bottom=2.2cm, left=2.5cm, right=2.5cm]{geometry}
%\usepackage{lmodern}
%\usepackage{microtype}
\usepackage{setspace}
\usepackage{titlesec}
\usepackage{titling}
\usepackage{abstract}
\usepackage{fancyhdr}
\usepackage{booktabs}
\usepackage{tabularx}
\usepackage{multirow}
\usepackage{array}
\usepackage{xcolor}
\usepackage{colortbl}
\usepackage{graphicx}
\usepackage{float}
\usepackage{caption}
\usepackage{subcaption}
\usepackage{amsmath}
\usepackage{amssymb}
\usepackage{enumitem}
\usepackage{hyperref}
\usepackage{tcolorbox}
\usepackage{mdframed}
\usepackage{parskip}

% ── Couleurs ──────────────────────────────────────────────────────────────────
\definecolor{bleuCMR}{RGB}{0, 62, 128}
\definecolor{vertCMR}{RGB}{0, 150, 57}
\definecolor{rougeCMR}{RGB}{206, 17, 38}
\definecolor{grisClair}{RGB}{245, 246, 248}
\definecolor{grisTexte}{RGB}{80, 80, 80}
\definecolor{jauneAlert}{RGB}{255, 243, 205}
\definecolor{jauneBord}{RGB}{230, 180, 0}

% ── Mise en page ──────────────────────────────────────────────────────────────
\onehalfspacing
\setlength{\parindent}{0pt}
\setlength{\parskip}{6pt}

% ── Titres de sections ────────────────────────────────────────────────────────
\titleformat{\section}{\large\bfseries\color{bleuCMR}}{\thesection.}{0.6em}{}[\color{bleuCMR}\titlerule]
\titleformat{\subsection}{\normalsize\bfseries\color{bleuCMR!80}}{\thesubsection.}{0.5em}{}
\titleformat{\subsubsection}{\normalsize\itshape\color{grisTexte}}{\thesubsubsection.}{0.5em}{}

\titlespacing*{\section}{0pt}{14pt}{6pt}
\titlespacing*{\subsection}{0pt}{10pt}{4pt}

% ── En-tête et pied de page ───────────────────────────────────────────────────
\pagestyle{fancy}
\fancyhf{}
\fancyhead[L]{\small\textcolor{grisTexte}{Audit sémantique LF Cameroun 2024--2025}}
\fancyhead[R]{\small\textcolor{grisTexte}{ISSEA -- ISE3 Data Science}}
\fancyfoot[C]{\small\thepage}
\renewcommand{\headrulewidth}{0.4pt}

% ── Encadré résultat clé ──────────────────────────────────────────────────────
\newmdenv[
  backgroundcolor=grisClair,
  linecolor=bleuCMR,
  linewidth=1.5pt,
  leftline=true, rightline=false, topline=false, bottomline=false,
  innerleftmargin=10pt, innerrightmargin=8pt,
  innertopmargin=6pt, innerbottommargin=6pt,
  skipabove=8pt, skipbelow=6pt
]{keybox}

\newmdenv[
  backgroundcolor=jauneAlert,
  linecolor=jauneBord,
  linewidth=1.5pt,
  leftline=true, rightline=false, topline=false, bottomline=false,
  innerleftmargin=10pt, innerrightmargin=8pt,
  innertopmargin=6pt, innerbottommargin=6pt,
  skipabove=6pt, skipbelow=4pt
]{alertbox}

% ── Commandes utiles ──────────────────────────────────────────────────────────
\newcommand{\mds}[1]{\textbf{#1~Mds FCFA}}
\newcommand{\pct}[1]{\textbf{#1\,\%}}
\newcommand{\sig}[1]{\textcolor{rougeCMR}{\textbf{#1}}}
\newcommand{\figref}[1]{\textcolor{bleuCMR}{\textit{[#1]}}}

% ── Méta-informations ────────────────────────────────────────────────────────
\hypersetup{
  colorlinks=true,
  linkcolor=bleuCMR,
  urlcolor=bleuCMR,
  pdftitle={Audit sémantique des Lois de Finances du Cameroun 2024-2025},
  pdfauthor={ISSEA -- ISE3 Data Science}
}

% ══════════════════════════════════════════════════════════════════════════════
\begin{document}
% ══════════════════════════════════════════════════════════════════════════════

% ── Page de titre ─────────────────────────────────────────────────────────────
\begin{titlepage}
\centering
\vspace*{1cm}

{\color{bleuCMR}\rule{\linewidth}{2pt}}
\vspace{0.4cm}

{\LARGE\bfseries\color{bleuCMR}
Audit Sémantique des Lois de Finances\\[0.3em]
du Cameroun\\[0.4em]
{\large\color{grisTexte} Exercices 2024 -- 2025}
}

\vspace{0.4cm}
{\color{bleuCMR}\rule{\linewidth}{2pt}}

\vspace{1.2cm}

{\large\textit{Analyse comparative par traitement automatique du langage,\\
modélisation thématique et classification zero-shot SND30}}

\vspace{1.5cm}

\begin{tabular}{ll}
\textbf{Institution} & Institut Sous-Régional de Statistique et\\
                     & d'Économie Appliquée (ISSEA), Yaoundé \\[4pt]
\textbf{Programme}   & ISE3 -- Spécialisation Data Science \\[4pt]
\textbf{Données}     & Extraction automatique via API Gemini \\[4pt]
                     & (Lois de Finances officielles du Cameroun) \\[4pt]
\textbf{Outils}      & Python 3.11, sentence-transformers,\\
                     & LDA, mDeBERTa-v3 (zero-shot), K-Means,\\
                     & HDBSCAN, scipy.stats \\[4pt]
\textbf{Date}        & Mars 2026 \\
\end{tabular}

\vfill

\begin{tcolorbox}[colback=grisClair, colframe=bleuCMR, width=0.9\linewidth,
  title={\textbf{Résumé exécutif}}, fonttitle=\small\bfseries]
\small
Ce rapport présente les résultats d'un audit sémantique comparatif des Lois de Finances (LF) du
Cameroun pour les exercices 2024 et 2025, portant sur \textbf{183 lignes budgétaires}, \textbf{256 articles} et 
\textbf{365 objectifs programmatiques}. L'analyse mobilise l'extraction automatique de texte via l'API Gemini, 
les embeddings multilingues (similarité cosinus), la modélisation LDA et la classification zero-shot alignée 
sur les 4 piliers de la Stratégie Nationale de Développement (SND30). 
Les résultats montrent un glissement sémantique faible mais statistiquement avéré 
(\mbox{Cohen's $d = -0{,}835$}, \mbox{$p < 0{,}001$}), 
une architecture programmatique stable (\pct{94{,}8} des programmes maintiennent leur pilier SND30), 
et un rééquilibrage budgétaire marqué en faveur du pilier Transformation structurelle (+\pct{143{,}5}).
\end{tcolorbox}

\end{titlepage}

\tableofcontents
\thispagestyle{fancy}
\newpage

% ══════════════════════════════════════════════════════════════════════════════
\section{Introduction}
% ══════════════════════════════════════════════════════════════════════════════

\subsection{Contexte}

La Loi de Finances (LF) est l'instrument juridique central de la politique budgétaire du Cameroun. 
Elle définit chaque année les ressources et les charges de l'État, en opérationnalisant les 
orientations de la Stratégie Nationale de Développement 2020--2030 (SND30) à travers quatre piliers 
stratégiques\,: la Transformation structurelle de l'économie, le Développement du capital humain, 
la Promotion de l'emploi et de l'insertion socio-économique, et la Gouvernance.

Dans un contexte de transition vers une budgétisation axée sur les résultats, évaluer la \textbf{continuité 
sémantique} entre les lois de finances successives constitue un enjeu analytique majeur. Un glissement 
lexical non maîtrisé peut signaler une rupture dans la cohérence programmatique, diluer l'efficacité 
des politiques publiques, ou révéler des réorientations stratégiques non explicitées.

\subsection{Problématiques}

Cette étude s'articule autour de trois questions de recherche\,:

\begin{enumerate}[label=\textbf{Q\arabic*.}, leftmargin=2cm]
  \item \textbf{Glissement sémantique\,:} Dans quelle mesure le contenu textuel de la LF 2025 
  diffère-t-il de celui de la LF 2024, et ce glissement est-il statistiquement significatif~?
  
  \item \textbf{Cohérence thématique\,:} La distribution thématique des articles et des objectifs 
  budgétaires a-t-elle évolué de manière significative entre les deux exercices~?
  
  \item \textbf{Alignement SND30\,:} L'architecture programmatique des objectifs budgétaires reste-t-elle 
  alignée sur les quatre piliers de la SND30~? Observe-t-on des reclassifications inter-piliers~?
\end{enumerate}

\subsection{Objectifs}

L'objectif général est de conduire un audit sémantique rigoureux et reproductible des LF 2024 et 2025 
du Cameroun, en mobilisant des méthodes de traitement automatique du langage (TAL) et d'analyse statistique.

Les objectifs spécifiques sont\,:
\begin{itemize}[noitemsep]
  \item Mesurer le glissement sémantique inter-annuel via la similarité cosinus d'embeddings multilingues~;
  \item Identifier l'évolution de la distribution thématique par modélisation LDA~;
  \item Évaluer l'alignement des objectifs budgétaires sur les piliers SND30 par classification zero-shot~;
  \item Quantifier la concentration budgétaire via les indices de Gini et HHI~;
  \item Tester statistiquement la significativité de chaque changement identifié.
\end{itemize}

% ══════════════════════════════════════════════════════════════════════════════
\section{Méthodologie}
% ══════════════════════════════════════════════════════════════════════════════

\subsection{Collecte et préparation des données}

Les textes des Lois de Finances 2024 et 2025 ont été extraits de manière automatique via l'\textbf{API 
Gemini} de Google (modèle \texttt{gemini-pro}), permettant le parsing structuré des documents PDF 
officiels publiés par le Ministère des Finances du Cameroun. Cette approche a permis d'obtenir un 
corpus segmenté en articles juridiques, chapitres et programmes budgétaires, avec les montants 
associés (Autorisations d'Engagement -- AE, et Crédits de Paiement -- CP).

\begin{keybox}
\small\textbf{Corpus final\,:} LF 2024\,: 168 articles, 183 lignes budgétaires, 183 objectifs~; 
LF 2025\,: 88 articles, 182 lignes budgétaires, 182 objectifs. Au total\,: 256 articles et 365 objectifs analysés.
\end{keybox}

\subsection{Pipeline analytique}

Le pipeline d'analyse comprend cinq étapes complémentaires, décrites ci-après et synthétisées 
dans la Figure~1.

\subsubsection{Embeddings et similarité cosinus}

Chaque article a été encodé via le modèle \texttt{paraphrase-multilingual-mpnet-base-v2} 
(sentence-transformers, 768 dimensions). Pour chaque article de la LF 2025, le \textbf{score 
de similarité cosinus maximal} avec l'ensemble des articles de la LF 2024 a été calculé\,:
\[
  s_i = \max_{j \in \text{LF}_{2024}} \frac{\mathbf{e}_{2025,i} \cdot \mathbf{e}_{2024,j}}
  {\|\mathbf{e}_{2025,i}\| \cdot \|\mathbf{e}_{2024,j}\|}
\]
Un seuil $\theta = 0{,}70$ a été fixé pour identifier les articles à changement sémantique significatif.

\subsubsection{Modélisation thématique LDA}

Un modèle LDA (\textit{Latent Dirichlet Allocation}) à 4 topics a été entraîné séparément sur 
les corpus 2024 et 2025, après nettoyage et lemmatisation. Les topics ont été interprétés comme\,: 
T0 -- Fiscalité \& revenus~; T1 -- Sanctions \& sécurité sociale~; T2 -- Administration fiscale~; 
T3 -- Finances publiques \& secteurs.

\subsubsection{Classification zero-shot SND30}

Chaque objectif budgétaire a été classifié selon les 4 piliers SND30 via le modèle 
\texttt{mDeBERTa-v3-base-mnli-xnli} (classification zero-shot multilingue). Un score de confiance 
$[0,1]$ est attribué à chaque pilier, et le pilier dominant est retenu.

\subsubsection{Clustering}

Deux algorithmes complémentaires ont été appliqués sur les représentations vectorielles des 
articles\,: \textbf{K-Means} ($k=3$, méthode du coude) et \textbf{HDBSCAN} (densité, 
$\texttt{min\_cluster\_size}=5$), permettant d'évaluer la structure latente du corpus.

\subsubsection{Tests statistiques}

Les différences inter-annuelles ont été testées par\,: le test de Kolmogorov-Smirnov (distributions), 
le test de Mann-Whitney U (comparaisons non-paramétriques), le test du Chi$^2$ (distributions catégorielles), 
le test de Kruskal-Wallis (comparaisons multi-groupes), et le $d$ de Cohen (taille d'effet).

% ══════════════════════════════════════════════════════════════════════════════
\section{Résultats}
% ══════════════════════════════════════════════════════════════════════════════

\subsection{Portrait macrobudgétaire 2024--2025}

\begin{table}[H]
\centering
\caption{Principaux indicateurs macrobudgétaires -- LF 2024 vs LF 2025}
\label{tab:macro}
\small
\begin{tabularx}{\linewidth}{lXXc}
\toprule
\rowcolor{bleuCMR!15}
\textbf{Indicateur} & \textbf{LF 2024} & \textbf{LF 2025} & \textbf{Variation} \\
\midrule
AE totales & 6,76 Mds FCFA & 7,38 Mds FCFA & \sig{+9,1\,\%} \\
CP totaux  & 6,68 Mds FCFA & 7,25 Mds FCFA & \sig{+8,6\,\%} \\
Ratio AE/CP & 1,013 & 1,018 & $\approx$ stable \\
Écart AE$-$CP & 85 M FCFA & 130 M FCFA & -- \\
Nb lignes budgétaires & 183 & 182 & $-1$ \\
Nb articles & 168 & 88 & $-47{,}6\,\%$ \\
Indice de Gini AE & 0,810 & 0,794 & $-0{,}016$ \\
HHI inter-piliers & 4\,427 & 4\,513 & $+85$ pts \\
\bottomrule
\end{tabularx}
\end{table}

Le budget 2025 affiche une croissance maîtrisée de \pct{9,1} en AE. La stabilité du ratio AE/CP 
à 1,01 traduit une gestion saine des engagements, sans accumulation de restes-à-payer. 
Le nombre de lignes budgétaires est quasi-identique (183$\to$182), confirmant la continuité 
de la nomenclature programmatique.

\figref{Insérer ici la Figure~1 -- Portrait macrobudgétaire\,: barres AE/CP, waterfall piliers, donut SND30}

\subsection{Évolution budgétaire par pilier SND30}

\begin{table}[H]
\centering
\caption{Évolution des AE par pilier SND30 (2024--2025)}
\label{tab:piliers}
\small
\begin{tabularx}{\linewidth}{lXXXXc}
\toprule
\rowcolor{bleuCMR!15}
\textbf{Pilier SND30} & \textbf{AE 2024} & \textbf{\% 2024} & \textbf{AE 2025} & \textbf{\% 2025} & \textbf{Évol.} \\
\midrule
Emploi et insertion & 4,06 Mds & 60,0\,\% & 4,57 Mds & 61,9\,\% & \sig{+12,6\,\%} \\
Gouvernance & 1,75 Mds & 25,9\,\% & 1,73 Mds & 23,4\,\% & $-1{,}3\,\%$ \\
Capital humain & 0,82 Mds & 12,2\,\% & 0,77 Mds & 10,4\,\% & \sig{$-6{,}5\,\%$} \\
Transform. structurelle & 0,13 Mds & 1,9\,\% & 0,31 Mds & 4,2\,\% & \sig{+143,5\,\%} \\
\bottomrule
\end{tabularx}
\end{table}

Deux signaux stratégiques se dégagent de cette comparaison\,: 

\textbf{(i) La Transformation structurelle triple sa dotation (+143,5\,\%)} passant de 127\,M à 309\,M FCFA. 
Bien que sa part absolue reste modeste (4,2\,\% des AE), cette progression marque une inflexion 
délibérée vers l'investissement productif, cohérente avec les ambitions industrielles de la SND30.

\textbf{(ii) Le Capital humain recule simultanément en nombre d'objectifs ($-11{,}4\,\%$, 35$\to$31) 
et en AE ($-6{,}5\,\%$, $-53$\,M FCFA)}, ramenant sa part de 12,2\,\% à 10,4\,\%. Ce double recul, 
portant sur l'éducation et la santé, mérite une surveillance particulière au regard des cibles SND30.

\begin{alertbox}
\small $\Rightarrow$~Le HHI inter-piliers augmente de 4\,427 à 4\,513 (niveau «~très concentré~»$>$2\,500), 
confirmant la polarisation croissante du budget vers le pilier Emploi et insertion ($\approx$\,62\,\% des AE).
\end{alertbox}

\subsection{Analyse du glissement sémantique}

\begin{table}[H]
\centering
\caption{Statistiques descriptives de la similarité cosinus par exercice}
\label{tab:sim}
\small
\begin{tabularx}{\linewidth}{lXXXXXXXc}
\toprule
\rowcolor{bleuCMR!15}
\textbf{Année} & \textbf{n} & \textbf{Moy.} & \textbf{Méd.} & \textbf{Éc-t.} & \textbf{Q1} & \textbf{Q3} & \textbf{Min} & \textbf{\%$<\theta$} \\
\midrule
2024 & 168 & 0,856 & 0,849 & 0,083 & 0,795 & 0,932 & 0,567 & 3,0\,\% \\
2025 & 88  & 0,916 & 0,931 & 0,060 & 0,879 & 0,961 & 0,699 & 1,1\,\% \\
\bottomrule
\end{tabularx}
\end{table}

Les scores de similarité cosinus sont élevés dans les deux exercices (0,856 et 0,916), confirmant 
une forte continuité textuelle. Seuls \textbf{5 articles sur 168 (3,0\,\%)} en 2024 et 
\textbf{1 article sur 88 (1,1\,\%)} en 2025 franchissent le seuil de rupture sémantique $\theta = 0{,}70$.

Trois tests statistiques convergent pour établir que les deux distributions sont structurellement distinctes\,:

\begin{table}[H]
\centering
\caption{Tests statistiques sur les distributions de similarité cosinus}
\label{tab:tests_sim}
\small
\begin{tabularx}{\linewidth}{lXXc}
\toprule
\rowcolor{bleuCMR!15}
\textbf{Test} & \textbf{Statistique} & \textbf{$p$-valeur} & \textbf{Conclusion} \\
\midrule
Kolmogorov-Smirnov & $D = 0{,}419$ & $<0{,}001$ & \sig{$\bigstar\bigstar\bigstar$} Sig. \\
Mann-Whitney U & $U = 4\,212$ & $<0{,}001$ & \sig{$\bigstar\bigstar\bigstar$} Sig. \\
Cohen's $d$ & $d = -0{,}835$ & -- & \sig{Effet fort} \\
Shapiro-Wilk 2024 & $W =\,--$ & $0{,}0008$ & Non-normal \\
Shapiro-Wilk 2025 & $W =\,--$ & $0{,}0002$ & Non-normal \\
\bottomrule
\end{tabularx}
\end{table}

Le rejet de la normalité dans les deux distributions (Shapiro-Wilk, $p<0{,}001$) justifie 
\textit{a posteriori} l'emploi de tests non-paramétriques. Le \textbf{Cohen's $d = -0{,}835$} 
constitue le résultat le plus informatif\,: la taille d'effet forte signifie que la loi 2025 
est \textit{structurellement plus proche} de la loi 2024 que l'inverse — le glissement est réel 
mais orienté vers plus de continuité, non vers la rupture.

\figref{Insérer ici la Figure~2\,: KDE + ECDF comparées, violin plots, profil de glissement par article}

\subsection{Analyse thématique par modélisation LDA}

La distribution des topics dominants révèle une restructuration thématique significative 
($\chi^2(3) = 8{,}14$, $p = 0{,}043$)\,:

\begin{table}[H]
\centering
\caption{Prévalence des topics dominants par exercice (LDA, 4 topics)}
\label{tab:topics}
\small
\begin{tabularx}{\linewidth}{lXXcc}
\toprule
\rowcolor{bleuCMR!15}
\textbf{Topic} & \textbf{2024 (n=168)} & \textbf{2025 (n=88)} & \textbf{$\Delta$ pts} & \textbf{Tendance} \\
\midrule
T0 – Fiscalité \& revenus      & 32 arts (19,0\,\%) &  8 arts (9,1\,\%) & $-9{,}9$ & $\downarrow$ \\
T1 – Sanctions \& séc. sociale & 33 arts (19,6\,\%) & 27 arts (30,7\,\%) & $+11{,}1$ & $\uparrow$ \\
T2 – Administration fiscale    & 61 arts (36,3\,\%) & 26 arts (29,5\,\%) & $-6{,}8$ & $\downarrow$ \\
T3 – Finances publ. \& sect.   & 42 arts (25,0\,\%) & 27 arts (30,7\,\%) & $+5{,}7$ & $\uparrow$ \\
\midrule
\multicolumn{4}{l}{\textit{Test Chi}$^2$(3) = 8,14 \quad $p$ = 0,043 \,\sig{$\bigstar$}} & \\
\bottomrule
\end{tabularx}
\end{table}

Le recul de \textbf{T0 (Fiscalité \& revenus)}, divisé par deux en proportion (19,0\,$\to$\,9,1\,\%), est 
le principal vecteur du changement. Il traduit une réduction des dispositions purement fiscales 
au profit de mécanismes réglementaires (T1, +11,1 pts) et d'une gestion sectorielle plus fine (T3, +5,7 pts).

L'analyse lexicale confirme ce glissement\,: en 2024, T0 est dominé par \textit{revenus, dépenses, 
exploitation} ; en 2025, il bascule vers \textit{taux, ministère, réalisation, mise en oe{}uvre, programmes} 
— transition du registre de la \textit{planification des recettes} vers celui de \textit{l'exécution 
programmatique}.

\figref{Insérer ici la Figure~3\,: prévalence comparée, top mots par topic, glissement lexical (delta probabilités)}

\subsection{Classification zero-shot et alignement SND30}

\subsubsection{Répartition des objectifs par pilier}

\begin{table}[H]
\centering
\caption{Objectifs budgétaires classifiés par pilier SND30}
\label{tab:snd30}
\small
\begin{tabularx}{\linewidth}{lXXXXXX}
\toprule
\rowcolor{bleuCMR!15}
\textbf{Pilier SND30} & \textbf{Obj. 2024} & \textbf{\% 2024} & \textbf{AE 2024} & \textbf{Obj. 2025} & \textbf{\% 2025} & \textbf{AE 2025} \\
\midrule
Emploi et insertion & 78 & 42,6\,\% & 4,06 Mds & 79 & 43,4\,\% & 4,57 Mds \\
Gouvernance & 65 & 35,5\,\% & 1,75 Mds & 66 & 36,3\,\% & 1,73 Mds \\
Capital humain & 35 & 19,1\,\% & 0,82 Mds & 31 & 17,0\,\% & 0,77 Mds \\
Transform. structurelle & 5 & 2,7\,\% & 0,13 Mds & 6 & 3,3\,\% & 0,31 Mds \\
\midrule
\multicolumn{2}{l}{\textit{Chi}$^2$(3)=0,34} & \multicolumn{4}{l}{$p$=0,951 — Distribution non significativement différente} & \\
\bottomrule
\end{tabularx}
\end{table}

Le Chi$^2$ quasi-nul ($p=0{,}951$) confirme que la \textbf{répartition catégorielle des objectifs 
par pilier est statistiquement identique entre 2024 et 2025}. Les tests de Mann-Whitney U sur les 
scores de classification confirment cette stabilité\,: tous les $p > 0{,}60$ et tous les 
Cohen's $d < 0{,}08$ (effet négligeable).

\subsubsection{Matrice de transition des piliers}

Sur \textbf{191 programmes communs} aux deux exercices, \textbf{181 (94,8\,\%)} maintiennent le même 
pilier SND30. Seulement 10 programmes (5,2\,\%) sont reclassifiés. Les flux principaux sont\,:

\begin{itemize}[noitemsep]
  \item \textbf{Emploi $\to$ Gouvernance (4 cas)\,:} \textit{Direction et coordination de l'action 
  gouvernementale}, \textit{Amélioration de la couverture sanitaire} -- révèlent la porosité sémantique 
  entre ces deux piliers sur les programmes transversaux.
  \item \textbf{Gouvernance $\to$ Emploi (3 cas)\,:} \textit{Gouvernance et appui institutionnel}, 
  \textit{Négociation et suivi de la coopération}.
  \item \textbf{Capital humain $\to$ Transformation structurelle (1 cas)\,:} \textit{Renforcement du 
  système de santé} -- interprétation cohérente avec une vision de la santé comme investissement économique.
\end{itemize}

\figref{Insérer ici la Figure~4\,: heatmap matrice de transition, scores de classification par pilier, heatmap scores moyens}

\subsection{Analyse du clustering}

\begin{table}[H]
\centering
\caption{Résultats du clustering K-Means et HDBSCAN}
\label{tab:clusters}
\small
\begin{tabularx}{\linewidth}{lXXXX}
\toprule
\rowcolor{bleuCMR!15}
\textbf{Algorithme} & \textbf{Exercice} & \textbf{Clusters} & \textbf{Distribution} & \textbf{Bruit} \\
\midrule
K-Means ($k=3$) & 2024 & 3 & C0:24 (14\,\%), C1:70 (42\,\%), C2:74 (44\,\%) & -- \\
K-Means ($k=3$) & 2025 & 3 & C0:28 (32\,\%), C1:20 (23\,\%), C2:40 (46\,\%) & -- \\
\midrule
HDBSCAN & 2024 & 4 & C0:7, C1:14, C2:40, C3:11 & 96 (57,1\,\%) \\
HDBSCAN & 2025 & 4 & C0:26, C1:9, C2:11, C3:5 & 37 (42,0\,\%) \\
\bottomrule
\end{tabularx}
\end{table}

Le taux de bruit HDBSCAN très élevé en 2024 (\textbf{57,1\,\%}) est le résultat le plus informatif\,: 
plus de la moitié des articles n'intègrent pas de clusters denses, confirmant l'\textbf{hétérogénéité 
intrinsèque du corpus juridique} -- un résultat attendu pour un texte couvrant de nombreux domaines. 
La réduction du bruit à 42,0\,\% en 2025 s'explique par la rationalisation du corpus ($-47{,}6\,\%$ d'articles).

La composition thématique des clusters évolue notablement\,: en 2024, les trois clusters sont 
thématiquement homogènes (tous dominés par T2 à $\approx$\,50$-$60\,\%). En 2025, une spécialisation 
émerge\,: \textbf{le Cluster 1 concentre 60,0\,\% d'articles T1} (Sanctions \& sécurité sociale), 
formant le groupe thématiquement le plus pur de l'ensemble du corpus. Le test de Kruskal-Wallis 
($H=5{,}20$, $p=0{,}074$) ne distingue pas les clusters par leur score de similarité cosinus.

\figref{Insérer ici la Figure~5\,: taille des clusters, taux de bruit HDBSCAN, composition thématique empilée}

\subsection{Concentration budgétaire}

Le Gini AE de \textbf{0,810 en 2024} et \textbf{0,794 en 2025} ($\Delta$ = $-0{,}016$) confirme une 
distribution extrêmement inégalitaire des dotations entre programmes\,: environ 20\,\% des programmes 
captent $\approx$\,80\,\% des ressources. La légère baisse ($-1{,}6$ point) indique une redistribution 
marginalement plus équitable, insuffisante pour modifier la structure de concentration.

Le diagramme de Pareto des AE 2025 révèle qu'un faible nombre de lignes budgétaires concentre 
l'essentiel des crédits, illustrant la loi de puissance caractéristique des budgets publics.

\begin{keybox}
\small Paradoxe Gini/HHI\,: le Gini intra-programme baisse légèrement ($-1{,}6$ pt) pendant que le HHI 
inter-piliers augmente (+85 pts). Ces deux mouvements simultanés reflètent une \textbf{consolidation 
stratégique}\,: les crédits sont marginalement mieux distribués au sein des programmes, mais la 
polarisation entre les 4 grands piliers s'accentue au profit de l'Emploi.
\end{keybox}

\figref{Insérer ici la Figure~6\,: courbes de Lorenz 2024 vs 2025, diagramme de Pareto AE 2025}

\subsection{Synthèse statistique}

\begin{table}[H]
\centering
\caption{Bilan de l'ensemble des tests statistiques réalisés}
\label{tab:bilan}
\small
\begin{tabularx}{\linewidth}{lXXc}
\toprule
\rowcolor{bleuCMR!15}
\textbf{Test} & \textbf{Objet} & \textbf{Statistique / $p$-val.} & \textbf{Décision} \\
\midrule
KS-test & Similarité cosinus 2024 vs 2025 & $D=0{,}419$ \,/\, $p<0{,}001$ & \sig{$\bigstar\bigstar\bigstar$ Sig.} \\
Mann-Whitney U & Similarité cosinus & $U=4\,212$ \,/\, $p<0{,}001$ & \sig{$\bigstar\bigstar\bigstar$ Sig.} \\
Chi$^2$ (topics) & Distrib. topics dominants & $\chi^2=8{,}14$ \,/\, $p=0{,}043$ & \sig{$\bigstar$ Sig.} \\
Chi$^2$ (piliers) & Distrib. piliers SND30 & $\chi^2=0{,}34$ \,/\, $p=0{,}951$ & n.s. \\
Mann-Whitney (piliers ×4) & Scores classif. SND30 & $d<0{,}08$ \,/\, $p>0{,}60$ & n.s. \\
Kruskal-Wallis & Similarité vs clusters & $H=5{,}20$ \,/\, $p=0{,}074$ & n.s. \\
\midrule
\multicolumn{3}{l}{\textbf{Bilan}\,: 3 tests significatifs sur 9 ($\alpha=5\,\%$)} & \\
\bottomrule
\end{tabularx}
\end{table}

Les tests significatifs portent sur le \textbf{glissement sémantique} (2 tests, effet fort) et la 
\textbf{distribution thématique LDA} (1 test, effet modéré). La stabilité programmatique SND30 est 
confirmée par 6 tests non-significatifs.

% ══════════════════════════════════════════════════════════════════════════════
\section{Recommandations}
% ══════════════════════════════════════════════════════════════════════════════

Sur la base des résultats présentés, six recommandations sont formulées à destination des décideurs 
et des équipes d'analyse budgétaire.

\begin{enumerate}[label=\textbf{R\arabic*.}, leftmargin=2cm, itemsep=6pt]

  \item \textbf{Renforcer le suivi du pilier Capital humain.}
  Le recul simultané du nombre d'objectifs ($35 \to 31$, $-11{,}4\,\%$) et des AE ($0{,}82 \to 0{,}77$ 
  Mds FCFA, $-6{,}5\,\%$) du pilier Capital humain doit faire l'objet d'un monitoring renforcé, en 
  particulier pour les sous-programmes Éducation et Santé, dont les effets de long terme sur la 
  croissance conditionnent l'atteinte des cibles SND30 à l'horizon 2030.

  \item \textbf{Valider manuellement les 10 reclassifications inter-piliers.}
  Les reclassifications de programmes d'un pilier SND30 à l'autre (notamment \textit{Couverture 
  sanitaire} d'Emploi vers Gouvernance, et \textit{Renforcement du système de santé} de Capital humain 
  vers Transformation structurelle) révèlent les limites du modèle zero-shot sur les intitulés 
  transversaux. Une revue manuelle de ces 10 cas est nécessaire pour garantir la robustesse de l'audit.

  \item \textbf{Appliquer la correction de Bonferroni pour les tests multiples.}
  Les 4 tests de Mann-Whitney U réalisés sur les scores de classification par pilier nécessitent 
  une correction du seuil ($\alpha_{\text{corr}} = 0{,}05/4 = 0{,}0125$) pour contrôler le risque 
  d'inflation du taux d'erreur de première espèce.

  \item \textbf{Optimiser le clustering HDBSCAN.}
  Le taux de bruit élevé (57,1\,\% en 2024) limite l'interprétabilité des clusters. L'exploration 
  de \texttt{min\_cluster\_size=3} combinée à une réduction dimensionnelle UMAP préalable devrait 
  permettre d'abaisser ce taux sous 30\,\% et de produire des groupes thématiques plus stables et 
  interprétables.

  \item \textbf{Croiser les AE avec les taux d'exécution réels.}
  L'audit porte sur les crédits \textit{programmés}. Son enrichissement avec les CP exécutés réels 
  permettrait de distinguer les ambitions budgétaires de la réalité de consommation des crédits, 
  dimension essentielle pour évaluer l'efficacité opérationnelle des politiques publiques.

  \item \textbf{Étendre l'analyse à la LF 2023 comme référence temporelle.}
  L'intégration d'un troisième exercice (LF 2023) permettrait de distinguer les tendances 
  structurelles pluriannuelles des variations conjoncturelles 2024--2025, et d'évaluer si le 
  rééquilibrage observé (hausse de la Transformation structurelle, recul du Capital humain) s'inscrit 
  dans une trajectoire continue ou constitue une rupture isolée.

\end{enumerate}

% ══════════════════════════════════════════════════════════════════════════════
\section*{Conclusion}
\addcontentsline{toc}{section}{Conclusion}
% ══════════════════════════════════════════════════════════════════════════════

Cet audit sémantique comparatif des Lois de Finances 2024 et 2025 du Cameroun apporte un éclairage 
quantitatif inédit sur la continuité et les évolutions de la programmation budgétaire nationale.

Les résultats établissent que la LF 2025 constitue une \textbf{évolution continue} de la LF 2024 
— non une rupture. La grande majorité des articles (97\,\%) présente une similarité cosinus supérieure 
au seuil critique $\theta=0{,}70$, et 94,8\,\% des programmes maintiennent leur alignement sur le 
même pilier SND30. Cette stabilité programmatique est confirmée par six tests statistiques non-significatifs.

Simultanément, deux mutations qualitatives méritent attention\,: (i) un \textbf{rééquilibrage 
budgétaire stratégique} vers la Transformation structurelle (+143,5\,\%) au détriment du Capital humain 
($-6{,}5\,\%$), et (ii) un \textbf{glissement thématique} des articles vers les dispositions réglementaires 
et la gestion sectorielle, statistiquement significatif ($\chi^2=8{,}14$, $p=0{,}043$).

La méthodologie mobilisée -- extraction automatique via API Gemini, embeddings multilingues, 
LDA, classification zero-shot et tests non-paramétriques -- constitue un cadre analytique reproductible 
et extensible aux exercices budgétaires futurs. Son intégration dans un tableau de bord dynamique 
de suivi de l'alignement SND30 représenterait un outil précieux pour le pilotage de la politique 
budgétaire camerounaise.

% ══════════════════════════════════════════════════════════════════════════════
\newpage
\appendix
\section*{Annexe -- Guide d'insertion des figures}
\addcontentsline{toc}{section}{Annexe -- Guide d'insertion des figures}
% ══════════════════════════════════════════════════════════════════════════════

\small
\begin{table}[H]
\centering
\caption*{Localisation et description des figures générées par le notebook}
\begin{tabularx}{\linewidth}{lXl}
\toprule
\rowcolor{bleuCMR!15}
\textbf{Figure} & \textbf{Description} & \textbf{Fichier} \\
\midrule
Fig.~1 & Portrait macrobudgétaire\,: barres AE/CP, waterfall piliers, donut SND30 & \texttt{fig1\_portrait\_budgetaire.png} \\
Fig.~2 & Glissement sémantique\,: KDE, ECDF, violin, profil similarité & \texttt{fig2\_glissement\_semantique.png} \\
Fig.~3 & Topic Modeling LDA\,: prévalence, top mots, glissement lexical & \texttt{fig3\_topic\_modeling.png} \\
Fig.~4 & Classification SND30\,: heatmap transition, scores, barres piliers & \texttt{fig4\_classification\_snd30.png} \\
Fig.~5 & Clustering\,: tailles K-Means, bruit HDBSCAN, composition thématique & \texttt{fig5\_clustering.png} \\
Fig.~6 & Concentration budgétaire\,: Lorenz, Pareto, courbe concentration & \texttt{fig6\_concentration.png} \\
Fig.~7 & Tableau de bord\,: volcano plot tests, Cohen's $d$, variations relatives & \texttt{fig7\_tableau\_bord\_stats.png} \\
\bottomrule
\end{tabularx}
\end{table}

\vspace{0.5cm}

\noindent\textbf{Instructions d'insertion dans ce document LaTeX\,:}

\begin{verbatim}
\begin{figure}[H]
  \centering
  \includegraphics[width=0.95\linewidth]{fig1_portrait_budgetaire.png}
  \caption{Portrait macrobudgétaire -- LF 2024 vs LF 2025}
  \label{fig:portrait}
\end{figure}
\end{verbatim}

\noindent Les fichiers \texttt{.png} sont générés lors de l'exécution du notebook 
\texttt{audit\_semantique\_cameroun\_2024\_2025.ipynb} et sauvegardés dans le répertoire 
\texttt{outputs/}.

\end{document}
